\chapter{Metodyka eksperymentów}
\label{chap:eksperymenty}

Celem eksperymentów jest porównanie wpływu zastosowanego algorytmu
uczenia przez wzmacnianie, reprezentacji stanu oraz funkcji nagrody
na jakość strategii uzyskiwanej w środowisku Ultimate Texas Hold'em.
W badaniu wykorzystano dwa algorytmy opisane w rozdziale
\ref{chap:metody-rl}: DQN oraz REINFORCE.

Eksperyment zaprojektowano w taki sposób, aby poszczególne warianty
różniły się wyłącznie badanymi elementami. Pozostałe właściwości
środowiska, sposób generowania rozdań oraz procedura ewaluacji są
wspólne dla wszystkich konfiguracji.

\section{Macierz eksperymentalna}

Porównanie obejmuje dwa algorytmy uczenia, dwie reprezentacje stanu
oraz dwie funkcje nagrody. Łącznie daje to osiem konfiguracji:

\begin{equation}
2 \times 2 \times 2 = 8.
\end{equation}

Dla każdego algorytmu analizowane są reprezentacje RAW oraz FEATURES.
Pierwsza z nich zawiera bezpośrednie kodowanie obserwowanych kart,
fazy rozdania oraz zbioru legalnych akcji. Druga rozszerza tę
reprezentację o cechy pokerowe wyznaczane na podstawie informacji
dostępnych graczowi.

W przypadku funkcji nagrody porównywana jest bezpośrednia wartość
zysku netto z rozdania oraz jej wersja przeskalowana przez maksymalną
wartość stawki. Skalowanie nie zmienia kierunku optymalizacji i nie
wprowadza obcinania wartości nagrody.

Ostateczna macierz obejmuje następujące konfiguracje:

\begin{table}[ht]
    \centering
    \caption{Konfiguracje wykorzystane w eksperymencie}
    \label{tab:experiment-configurations}
    \begin{tabular}{lll}
        \hline
        Algorytm & Reprezentacja & Nagroda \\
        \hline
        DQN & RAW & zysk netto \\
        DQN & RAW & zysk netto skalowany \\
        DQN & FEATURES & zysk netto \\
        DQN & FEATURES & zysk netto skalowany \\
        REINFORCE & RAW & zysk netto \\
        REINFORCE & RAW & zysk netto skalowany \\
        REINFORCE & FEATURES & zysk netto \\
        REINFORCE & FEATURES & zysk netto skalowany \\
        \hline
    \end{tabular}
\end{table}

\section{Powtarzalność eksperymentów}

Ze względu na losowy charakter procesu uczenia pojedynczy przebieg
nie jest wystarczający do wiarygodnej oceny konfiguracji.
Każdy z ośmiu wariantów jest dlatego trenowany z użyciem pięciu
niezależnych ziaren losowości. Prowadzi to do wykonania łącznie

\begin{equation}
8 \times 5 = 40
\end{equation}

niezależnych procesów treningowych.

Ten sam zestaw pięciu ziaren jest wykorzystywany dla wszystkich
wariantów. Takie rozwiązanie pozwala ograniczyć wpływ przypadkowych
różnic wynikających z losowania rozdań i ułatwia porównanie wpływu
reprezentacji stanu oraz funkcji nagrody.

Oddzielne generatory liczb pseudolosowych wykorzystywane są przez
środowisko oraz mechanizmy eksploracji agentów. Ewaluacja okresowa
nie modyfikuje stanu generatorów używanych podczas treningu, dzięki
czemu jej wykonywanie nie wpływa na dalszy przebieg uczenia.

\section{Budżet treningowy}

Podstawowy budżet treningowy ustalono na 50~000 epizodów dla każdej
konfiguracji. Wartość tę wybrano na podstawie wcześniejszego
eksperymentu diagnostycznego, w którym obserwowano zachowanie
wszystkich ośmiu wariantów podczas treningu.

W przypadku DQN zwiększanie liczby epizodów nie prowadziło do
systematycznego wzrostu jakości polityki. W części konfiguracji
obserwowano również pogorszenie wyniku po dłuższym treningu.

W przypadku REINFORCE część przebiegów osiągała stabilne, proste
polityki, takie jak wybieranie tego samego zakładu dla większości
stanów. Dodatkowy przebieg diagnostyczny rozszerzony do 100~000
epizodów wykazał, że dalsze zwiększanie budżetu nie prowadziło do
trwałej poprawy. Polityka przechodziła pomiędzy różnymi strategiami,
a następnie ponownie zbliżała się do wcześniej obserwowanego
rozwiązania.

Z tego powodu przyjęto wspólny budżet 50~000 epizodów dla obu
algorytmów. Pozwala to również zachować porównywalność kosztu
uczenia pomiędzy DQN i REINFORCE.

\section{Walidacja okresowa}

W trakcie treningu wykonywana jest okresowa ewaluacja aktualnej
polityki. Jej celem jest obserwacja procesu uczenia i budowa krzywych
przedstawiających zależność jakości strategii od liczby wykonanych
epizodów.

Walidacja wykonywana jest co 4~000 epizodów oraz dodatkowo po
zakończeniu 50~000 epizodów. Każdy punkt walidacyjny obejmuje
2~000 rozdań wygenerowanych na podstawie tego samego, ustalonego
harmonogramu talii.

Zastosowanie wspólnego harmonogramu oznacza, że kolejne wersje
polityki są oceniane na identycznym zestawie sytuacji pokerowych.
Zmiana wartości oczekiwanej pomiędzy punktami walidacyjnymi wynika
więc ze zmiany zachowania agenta, a nie z wykorzystania innego
zestawu kart.

Podczas walidacji wyłączona jest eksploracja. Agent DQN wybiera
legalną akcję o najwyższej wartości Q. W przypadku REINFORCE
wybierana jest legalna akcja o najwyższym prawdopodobieństwie
zgodnie z aktualną polityką. Pozwala to oceniać deterministyczną
strategię wynikającą z aktualnego stanu modelu.

Harmonogram używany podczas walidacji jest niezależny od rozdań
wykorzystywanych w procesie treningowym.

\section{Rozszerzona analiza treningu}

Po wykonaniu podstawowej macierzy eksperymentalnej planowana jest
dodatkowa analiza długiego treningu. Na podstawie wyników
walidacyjnych wybrany zostanie jeden wariant DQN oraz jeden wariant
REINFORCE.

Wybrane konfiguracje zostaną dodatkowo wytrenowane do 100~000
epizodów. Celem tej części badania nie jest zastąpienie podstawowego
budżetu 50~000 epizodów ani wyszukiwanie najlepszego chwilowego
checkpointu. Analiza ma umożliwić ocenę, czy dalszy trening prowadzi
do trwałej poprawy strategii, stabilizacji polityki lub jej dalszych
oscylacji.

Wyniki tej części są traktowane jako analiza uzupełniająca i nie
zmieniają podstawowej macierzy 40 procesów treningowych.

\section{Ewaluacja końcowa}

Po zakończeniu treningu modele są oceniane na oddzielnym,
nieużywanym wcześniej harmonogramie rozdań. Harmonogram końcowy
nie jest wykorzystywany podczas treningu, walidacji ani wyboru
budżetu treningowego.

Końcowa ewaluacja obejmuje 100~000 identycznych rozdań dla każdej
ocenianej strategii. Na tym samym zbiorze oceniani są również
RandomAgent oraz RuleBasedAgent. Agenci bazowi nie wymagają procesu
treningowego, dlatego liczba 100~000 dotyczy w ich przypadku
wyłącznie liczby rozegranych rund ewaluacyjnych.

Podstawową miarą jakości strategii jest średni zysk netto na jedno
rozdanie:

\begin{equation}
EV =
\frac{1}{N}
\sum_{i=1}^{N} p_i,
\end{equation}

gdzie $p_i$ oznacza zysk netto uzyskany w $i$-tym rozdaniu, natomiast
$N$ oznacza liczbę rozdań ewaluacyjnych.

Oprócz wartości oczekiwanej rejestrowane są odchylenie standardowe
wyników, błąd standardowy średniej, średnia wartość postawionych
żetonów, liczności wyników rozdania oraz częstości wyboru
poszczególnych akcji.

\section{Porównania sparowane}

Wszystkie strategie podczas ewaluacji końcowej rozgrywają ten sam
uporządkowany zestaw rozdań. Umożliwia to wykorzystanie porównań
sparowanych.

Dla dwóch porównywanych strategii $A$ oraz $B$ dla każdego rozdania
wyznaczana jest różnica

\begin{equation}
d_i = p_{A,i} - p_{B,i}.
\end{equation}

Średnia różnica wynosi

\begin{equation}
\bar{d} =
\frac{1}{N}
\sum_{i=1}^{N} d_i.
\end{equation}

Na podstawie rozkładu różnic wyznaczane są również odchylenie
standardowe, błąd standardowy oraz 95-procentowy przedział ufności.
Podejście sparowane ogranicza wpływ zmienności wynikającej z
różnych układów kart, ponieważ obie strategie oceniane są w tych
samych sytuacjach.

Analiza końcowa obejmuje porównanie algorytmów DQN i REINFORCE,
wpływu reprezentacji RAW i FEATURES, wpływu zastosowanej funkcji
nagrody oraz odniesienie uzyskanych wyników do agentów bazowych.